\section{Introduction}
Source code change history is valuable to both software practitioners and researchers. Practitioners leverage source code history for tasks such as code reviews, expert identification, and developer onboarding~\cite{grund_codeshovel_2021}. Researchers, on the other hand, use code history for a variety of tasks, including bug prediction~\cite{chowdhury_method-level_2024, mashhadi_empirical_2024}, code clone management~\cite{duala-ekoko_tracking_2007}, entity name recommendation~\cite{kashiwabara_method_2015}, security vulnerability detection \cite{li_cleanvul_2024}, code review \cite{grund_codeshovel_2021}, authorship attribution ~\cite{meng_mining_2013, rahman_ownership_2011}, and software maintenance effort estimation~\cite{chowdhury_good_2025, chowdhury_evidence_2025, hassan_versioned_2024}. To extract this historical information, researchers primarily relied on version control systems (VCS). However, traditional \texttt{VCSs} like \texttt{Git} mainly focus on file-level change tracking, which lacks the granularity needed for studies at finer levels. In practice, understanding maintenance activities at lower granularities—particularly at the method level—is important to both researchers and practitioners~\cite{grund_codeshovel_2021, chowdhury_evidence_2025, chowdhury_method-level_2024}. However, tracking method-level change history is challenging, as methods may be renamed or moved across files during refactoring and other development activities. These challenges and practical demands have driven the development of specialized tools for generating method-level change histories~\cite{higo_tracking_2020, grund_codeshovel_2021, jodavi_accurate_2022}.

Several approaches have been proposed to recover method-level code evolution, particularly from \texttt{Git} repositories. Hata \textit{et al.}~\cite{hata_historage_2011} introduced \textit{Historage}, a tool that enables method-level change tracking by restructuring repositories to store individual methods in separate files. This approach was later refined in \textit{FinerGit}~\cite{higo_tracking_2020}, which enhanced the accuracy of detecting small method changes. However, both tools require extensive preprocessing~\cite{grund_codeshovel_2021} of the codebase before any method-level history can be generated. This introduces additional storage overhead and preprocessing time, making them impractical for scenarios where fast and on-demand change history extraction is needed.

To mitigate the need for upfront preprocessing, Grund \textit{et al.} ~\cite{grund_codeshovel_2021} introduced \textit{CodeShovel} (ICSE 2021), a tool that tracks method-level change history on demand using a lightweight string similarity algorithm. \textit{CodeShovel} was reported to be both accurate and efficient, based on evaluation using a human-generated change history oracle of 200 Java methods. Unfortunately, Jodavi \textit{et al.}~\cite{jodavi_accurate_2022} later identified inaccuracies in this oracle. They corrected the oracle and introduced a more robust and refactoring-aware tool,  \textit{CodeTracker} (ESEC/FSE 2022), which significantly outperformed \textit{CodeShovel} in terms of precision and recall.

A concern remains, however, that the internal threats present in the human-generated \textit{CodeShovel} oracle may have carried over into the \textit{CodeTracker} oracle, which was largely derived from \textit{CodeShovel} and later corrected by a separate group of researchers, thereby introducing its own set of internal threats. Such inaccuracies in recorded change histories are particularly problematic, as tools validated by these oracles form the foundation for research on method-level maintenance analysis and modeling~\cite{ahmad_impact_2025, chowdhury_evidence_2025, chowdhury_good_2025, chowdhury_method-level_2024, chowdhury_revisiting_2022}. Moreover, the accuracy of tools evaluated against these less reliable oracles remains uncertain, raising concerns about their practical usefulness to software practitioners. Low precision in identifying method change commits can result in excessive false positives, while low recall may lead to missing critical changes—both of which can harm the usefulness of change history tracking in real-world scenarios.

To address the limitations of human-generated change histories, a more systematic, semi-automated approach was adopted to construct two new oracles for method-level change history. Rather than relying solely on manual annotation, A tool was developed to collect method-change commits for a given Java method using multiple state-of-the-art tools—namely, \textit{CodeShovel}, \textit{CodeTracker}, \textit{GitLog} (\textit{GitFuncName} and \textit{GitLineRange}), and \textit{IntelliJ}. The tool aggregates method change commits collected by all sources, based on the observation that this combined approach achieves significantly higher recall than any individual tool. This union set was then manually verified by two expert software engineers, who removed false positives to ensure high precision and added any relevant method change commits missed by all tools. The result is two high-quality oracles that combine the strengths of automated detection and expert validation. As shown later in the analysis (\textbf{RQ1}), the previously used oracles from \textit{CodeShovel} and \textit{CodeTracker} contain notable inaccuracies—confirming concerns about the validity of prior evaluations that relied on those oracles.

In addition to providing the new oracles, a novel, heuristic-based, method-level history generation algorithm and tool called \textit{HistoryFinder} is introduced. It also addresses several of \textit{CodeShovel}'s technical limitations. For instance, unlike \textit{CodeShovel}, \textit{HistoryFinder} captures changes introduced in merge commits, as well as modifications to JavaDocs, annotations, and code formatting.

\textit{HistoryFinder} is then compared with other method-level change-history tracking tools in terms of precision, recall, and runtime using the two newly developed oracles (\textbf{RQ2} and \textbf{RQ3}). Overall, the research-oriented tools—\textit{HistoryFinder}, \textit{CodeShovel}, and \textit{CodeTracker}—consistently demonstrate substantially higher accuracy than other alternatives such as \textit{GitLog} and \textit{IntelliJ}. \textit{HistoryFinder} generally achieves the highest precision and recall across all test oracles, while also offering competitive runtime performance. Although \textit{GitLog} is the fastest at reconstructing method histories, it suffers from significantly lower precision and recall, making \textit{HistoryFinder} the most balanced option in terms of accuracy and efficiency. To support broader adoption, a \texttt{GUI-based} tool is also provided that allows researchers and practitioners to generate method-level change histories using any of the state-of-the-art tools.

The contribution of this chapter can be summarized as follows.
\begin{itemize}
    \item It offers two systematically developed method-level change history oracles---consisting of 400 Java methods. Building these oracles with manual verification required approximately 300 hours from each of the two participating authors, excluding the time for building the necessary tools. These more accurate new oracles can be instrumental for future research on developing more method-level change history generation tools.
    \item It introduces a new algorithm and tool, \textit{HistoryFinder}, that addresses the limitations of the \textit{CodeShovel} tool.
    \item It conducts a thorough evaluation of the state-of-the-art method tracking tools against two accurate oracles to assess their relative performance in terms of precision, recall, F1 score, and runtime efficiency.
    \item It provides a user-friendly GUI-based tool that integrates all the state-of-the-art method-tracking tools (excluding \textit{IntelliJ}), allowing users to analyze and generate method histories for any arbitrary Java method from a \texttt{Git-based} project.
\end{itemize}

The two oracles and the GUI-based history generation tool are publicly shared.\footnote{\textit{HistoryFinder}: \url{https://github.com/SQMLab/HistoryFinder}} While this \texttt{GUI-based} tool can be used by both practitioners and researchers, the shared repository also shows how to use the new \textit{HistoryFinder} tool from the command line for enabling research in \emph{mining software repositories}.
